\scanpage{221}
\mspage{76}

$R$ sur $P\times Q$. Voici les conditions sur $\beta_0:P\to Q$ :

\textcircled{a}\quad $\beta_0$ est \uline{\unsure{décroissante|décroissant}},

\deleted{b) \unsure{pour} $\forall$ partie} \unsure{ce} qui implique que la relation

\[
 \underset{\substack{\cap\\P\times Q}}{R_{\beta_0}}\deleted{\subset}
 =\{(p,q)\in P\times Q\mid q\leq\beta_0(p)\}
\]

satisfait à la condition $(D_0)$ du côté gauche, et il faut encore la vérifier du côté droit, i.e regarder, pour $q$ fixé, l'ensemble

\[
 \{q\}^{\unsure{\cap|\frown}}=\{p\in P\mid q\leq\beta_0(p)\}.
\]

Comme $\beta_0$ est \unsure{décroissante|décroissant}, on voit que \deleted{\unsure{est}} $p\in\{q\}^\cap$, et $p'\leq p\Longrightarrow p'\in\{q\}^\cap$, i.e. $\{q\}^\cap$ est \uline{ouvert}. \deleted{Pour} On doit vérifier encore (supposant l'existence des Sup. quelconques dans $P$) qu'il \deleted{est} cet ensemble est stable par \deleted{\unsure{lim}} Sup. quelconques - il en résultera qu'il est de la forme $P_{/\alpha(q)}$. Cela équivaut à la condition

\textcircled{b}\quad $\beta_0$ transforme Sup quelconques \inserted{dans $P$} en Inf \inserted{dans $Q$}.

\deleted{\unsure{(}} Si on regarde $\beta_0$ comme un \inserted{\unsure{contr.|covar.}} foncteur $P^{\unsure{\circ|o}}\longrightarrow Q$, cela signifie qu'il commute aux limites projectives, ou encore qu'il a un adjoint à gauche $Q\longrightarrow P^\circ$, lequel commute aux limites inductives, et correspond
% archive sheet number: 221
